% 00_example.tex — образец. Новый документ: скопируйте этот файл в docs/0N_<имя>.tex,
% замените данные документа и пишите разделы по оглавлению из README.md §4.
% Собрать PDF: ./build.sh docs/00_example.tex → pdf/00_example.pdf
%
% Как писать текст:
%   - абзацы разделяются пустой строкой;
%   - «_» в ID пишется как есть: SYS_04_02;
%   - символы # % & $ { } экранируются: \# \% \& \$ \{ \};
%   - кавычки «ёлочки», тире — длинное «—», диапазон — «--» (стр. 75--76).
\documentclass{tpks}

% --- Данные документа (проект, команда и курс подставятся из project.tex) ---
\doctitle{Пример документа}
\docversion{1}
\docdate{02.10.2026}
% \docauthors{Иванов И. И., Петров П. П.}  % необязательно: по умолчанию — вся команда

\begin{document}
\makeheader % шапка и содержание; основной текст начинается со следующей страницы

\section{Введение}

\subsection{Назначение}

Этот файл показывает все элементы шаблона. Шапка, содержание, нумерация разделов, таблиц
и рисунков, лист регистрации изменений получаются сами.

\subsection{Термины и сокращения}

Определения берутся из \texttt{docs/glossary.tex}. В документе перечисляем только нужные:

\printterms{ЕК, СП, ЗИ, Базовая версия}

\section{Пример требований}

Одно требование — одно окружение \texttt{req}: первый аргумент — ID по схеме из
README.md §4.6, второй — версия. Повтор ID или ID не по схеме остановит сборку.

\begin{req}{SYS_04_01}{1}
Система должна обеспечивать продвижение релиза на следующий стенд только после прохождения
всех КТК текущего стенда.
\reqsource{Функция 04 «Продвижение релиза».}
\reqverify{Испытание: релиз с непройденной КТК не продвигается.}
\end{req}

\begin{req}{SYS_04_02}{2}
Система должна регистрировать в журнале каждое продвижение релиза с указанием инициатора,
времени и результата КТК.
\reqsource{Функция 04 «Продвижение релиза»; уточнено по замечанию от 17.10.2026.}
\reqverify{Испытание: после продвижения в журнале есть запись со всеми полями.}
\reqnote{Состав полей журнала задаётся в требованиях к ПО.}
\end{req}

Ссылка на требование в тексте: \reqref{SYS_04_01}. В этом документе она кликабельна,
ссылка на требование из другого документа печатается просто как ID: \reqref{SYS_01_03}.

Вопрос, который нельзя решить без людей, помечаем так: \clarify{кто подтверждает выход
на рабочий стенд}. Пустая пометка: \clarify{}.

\section{Пример таблицы и списков}

Таблица~\ref{tab:states} перечисляет состояния релиза. Колонка \texttt{L} занимает
оставшуюся ширину и переносит текст, колонка \texttt{l} — по ширине содержимого.
Последний аргумент — заголовок столбцов: он повторится, если таблица перейдёт на новую страницу.

\begin{tabl}{|l|L|}{Состояния релиза}[tab:states]{\textbf{Состояние} & \textbf{Что значит}}
Черновик & Релиз создан, состав ещё меняется \\ \hline
Собран & Сборка прошла, артефакты лежат в хранилище дистрибутивов \\ \hline
Выпущен & Релиз установлен на рабочий стенд \\ \hline
\end{tabl}

Маркированный список:
\begin{itemize}
  \item первый пункт;
  \item второй пункт.
\end{itemize}

Нумерованный список:
\begin{enumerate}
  \item Зарегистрировать СП.
  \item Получить решение совета.
\end{enumerate}

% Рисунок: положите файл в docs/img/ и раскомментируйте.
% \begin{figure}[h]
%   \centering
%   \includegraphics[width=0.8\textwidth]{release_states.png}
%   \caption{Жизненный цикл релиза}\label{fig:states}
% \end{figure}

\section{Команда}

Состав подставляется из \texttt{docs/project.tex}:

\printteam

% Лист регистрации изменений — всегда в конце документа. Новая версия — новая строка.
\begin{changelog}
\change{1}{02.10.2026}{Первая версия}{Брылкина К. К.}{—}
\change{2}{18.10.2026}{Уточнено требование SYS_04_02}{Брылкина К. К.}{Замечание преподавателя от 17.10.2026}
\end{changelog}

\end{document}
